\documentclass[a4paper]{article}

% Basic packages
% Español (México) con XeLaTeX: fontspec para fuentes Unicode, polyglossia
% para la división silábica y los nombres del documento en la variante mexicana.
\usepackage{fontspec}
\usepackage{polyglossia}
\setdefaultlanguage[variant=mexican]{spanish}
% Los nombres chinos van como «pinyin (original)»: xeCJK da a los caracteres
% Han su propia fuente; sin ella XeLaTeX los omite con un aviso.
% FandolSong no cubre algunos caracteres de nombres (U+73FA); WenQuanYi Zen Hei sí.
\usepackage[AutoFallBack=true]{xeCJK}
\setCJKmainfont{FandolSong-Regular.otf}
\setCJKfallbackfamilyfont{\CJKrmdefault}{WenQuanYi Zen Hei}
% polyglossia no traduce los nombres de listings.
\AtBeginDocument{\providecommand{\lstlistingname}{}\renewcommand{\lstlistingname}{Listado}%
  \providecommand{\lstlistlistingname}{}\renewcommand{\lstlistlistingname}{Índice de listados}}
\usepackage{amsmath, amssymb}
\usepackage{graphicx}
\usepackage[margin=2.5cm]{geometry}
\usepackage[most]{tcolorbox}
\usepackage{etoolbox}
\usepackage{listings}
\usepackage{hyperref}
\usepackage{booktabs}
\usepackage{subcaption}
\usepackage{float}

% Highlight boxes
\newtcolorbox{knowledgebox}[1]{
    enhanced, colback=blue!5!white, colframe=blue!75!black, colbacktitle=blue!75!black,
    coltitle=white, fonttitle=\bfseries, title={#1}, attach boxed title to top left={yshift=-2mm, xshift=2mm},
    boxrule=1pt, sharp corners
}

\newtcolorbox{importantbox}[1]{
    enhanced, colback=yellow!10!white, colframe=yellow!80!black, colbacktitle=yellow!80!black,
    coltitle=black, fonttitle=\bfseries, title={#1}, sharp corners
}

\newtcolorbox{warningbox}[1]{
    enhanced, colback=red!5!white, colframe=red!75!black, colbacktitle=red!75!black,
    coltitle=white, fonttitle=\bfseries, title={#1}, sharp corners
}

% Listings
\lstset{
    language=Python,
    basicstyle=\ttfamily\small,
    keywordstyle=\color{blue},
    stringstyle=\color{red!60!black},
    commentstyle=\color{green!60!black},
    breaklines=true,
    frame=single,
    numbers=left,
    numberstyle=\tiny\color{gray},
    captionpos=b,
    extendedchars=true
}

% Front-page metadata
\newcommand{\notetitle}{CS146S Week 4: Coding Agent Patterns}
\newcommand{\noteauthors}{Elaboradas a partir de la clase de Mihail Eric + Boris Cherny (Claude Code)}
\newcommand{\notedate}{Octubre de 2025}
\newcommand{\videochannel}{Stanford CS146S}
\newcommand{\videopublishdate}{Fall 2025}
\newcommand{\videoduration}{2 lectures (Mon + Fri)}
\newcommand{\videourl}{https://themodernsoftware.dev}
\newcommand{\videocoverpath}{}
\newcommand{\repourl}{https://github.com/hqhq1025/ai-course-notes}


% Footer with repo info
\usepackage{fancyhdr}
\pagestyle{fancy}
\fancyhf{}
\fancyfoot[C]{\thepage}
\fancyfoot[R]{\footnotesize\color{gray}\href{\repourl}{AI Course Notes}}
\renewcommand{\headrulewidth}{0pt}
\renewcommand{\footrulewidth}{0pt}

\begin{document}

\begin{titlepage}
\centering
{\Large Notas del curso\par}
\vspace{1.2cm}
{\huge\bfseries \notetitle\par}
\vspace{0.8cm}
{\large \noteauthors\par}
\vspace{0.3cm}
{\large \notedate\par}
\vspace{1.2cm}

\ifdefempty{\videocoverpath}{
\vspace{3cm}
}{
\includegraphics[width=0.82\textwidth,height=0.45\textheight,keepaspectratio]{\videocoverpath}\par
}

\vfill
\begin{tcolorbox}[width=0.9\textwidth, colback=black!2!white, colframe=black!60, sharp corners]
\textbf{Curso}: \videochannel\par
\textbf{Semestre}: \videopublishdate\par
\textbf{Sesión}: \videoduration\par
\textbf{Sitio}: \href{\videourl}{\nolinkurl{\videourl}}\par
\textbf{Página del proyecto}: \href{\repourl}{\nolinkurl{\repourl}}\par
\end{tcolorbox}
\end{titlepage}

\tableofcontents
\newpage

%% --- Inicio del contenido --- %%

\section{Cómo convertirse en Agent Manager}

\subsection{Evolución de los equipos de desarrollo de software}

La forma de organización de los equipos de desarrollo de software está atravesando su cuarta transformación importante:

\begin{enumerate}
    \item \textbf{1940--1960s}: un solo desarrollador gestiona todo el proyecto
    \item \textbf{1970--1990s}: los equipos de software se vuelven la norma y surge la división especializada del trabajo
    \item \textbf{2023--2025}: los desarrolladores se apoyan en sistemas de programación con IA para su trabajo
    \item \textbf{2025+}: un solo desarrollador gestiona la producción de trabajo de múltiples AI Agents
\end{enumerate}

\begin{importantbox}{De gestionar personas a gestionar Agents}
La ruta de evolución del desarrollo es: gestionar la propia producción $\rightarrow$ gestionar la producción del equipo $\rightarrow$ gestionar la producción del equipo asistido por IA $\rightarrow$ \textbf{una sola persona gestiona la producción de múltiples AI Agents}. Este último paso es la transformación que está ocurriendo en este momento.
\end{importantbox}

\subsection{Descomposición de las tareas de software}

Una tarea típica de desarrollo de software puede descomponerse en los siguientes pasos, donde el verde indica que la dirige un humano y el azul indica que puede encargarse un Agent:

\begin{enumerate}
    \item Proporcionar los requisitos de alto nivel \textbf{{[}Humano{]}}
    \item Convertir los requisitos en un documento de diseño \textbf{{[}Humano/Agent{]}}
    \item Implementar la solución según el documento de diseño \textbf{{[}Agent{]}}
    \item Agregar pruebas \textbf{{[}Agent{]}}
    \item Asegurar que el CI pase \textbf{{[}Agent{]}}
    \item Revisión de código \textbf{{[}Agent{]}}
    \item Actualizar la documentación \textbf{{[}Agent{]}}
\end{enumerate}

\begin{knowledgebox}{El valor central del humano está en los pasos 1--2}
En el modelo de Agent Manager, el mayor valor del ingeniero humano se manifiesta en la \textbf{comprensión de los requisitos} y el \textbf{diseño de la arquitectura}. Estos dos pasos requieren contexto de negocio, juicio técnico y «gusto»: precisamente el eslabón más débil de la IA actual.
\end{knowledgebox}

\subsection{Cuatro técnicas para guiar a un Agent}

\subsubsection{Archivos de comportamiento del Agent}

Como CLAUDE.md, .cursorrules, AGENTS.md, entre otros, que ofrecen una guía de comportamiento a nivel de proyecto.

\subsubsection{Hooks}

Scripts deterministas que se ejecutan cuando se disparan eventos predefinidos:
\begin{itemize}
    \item \texttt{PreToolUse}: antes de invocar una herramienta
    \item \texttt{PostToolUse}: después de invocar una herramienta
    \item \texttt{UserPromptSubmit}: cuando el usuario envía un prompt
    \item \texttt{PreCompact}: antes de comprimir el context
\end{itemize}

\subsubsection{Commands}

Guardar los prompts de uso frecuente como archivos que el Agent puede ejecutar según se necesite. Casos de uso típicos: ejecutar pruebas, revisar código, hacer un commit de git y subir los cambios.

\subsubsection{Subagents}

Mecanismo de delegación en tiempo de ejecución, con los siguientes objetivos centrales:
\begin{itemize}
    \item Crear «personalidades» de desarrollador independientes para distintos tipos de trabajo (frontend, backend, etc.)
    \item Separar con claridad el context de distintos flujos de trabajo
    \item Proporcionar un system prompt personalizado, un conjunto de herramientas y una context window independiente
\end{itemize}

\begin{importantbox}{El Subagent es el embrión de un Agent que gestiona a otro Agent}
El patrón de Subagent representa una tendencia importante: un Agent gestiona a otros Agents. El Agent principal se encarga de descomponer y delegar tareas, el subagent se concentra en la ejecución de un dominio específico, y cada subagent cuenta con un context independiente para evitar la sobrecarga de información.
\end{importantbox}

\subsection{Mejores prácticas}

\begin{enumerate}
    \item \textbf{Establecer barreras de protección estrictas}: las pruebas del código base más las mejores prácticas de CI/CD son un backstop indispensable
    \item \textbf{Auditar las modificaciones de cada Agent}: etiquetar (label) cada diff generado por un Agent
    \item \textbf{Usar distintos modelos según la tarea}: modelos rápidos para tareas simples, modelos potentes para tareas complejas
    \item \textbf{Las tareas complejas requieren más orientación previa}: no todas las tareas son adecuadas para una delegación completamente asíncrona
    \item \textbf{Hacer checkpoint (commit) con frecuencia}: hacer commit del código de manera periódica, para evitar que se pierda el avance si el Agent se desvía de rumbo
\end{enumerate}

\begin{warningbox}{Preguntas abiertas}
\begin{itemize}
    \item ¿Cómo automatizar el 10--20\% inicial de investigación de cada tarea?
    \item ¿Cómo mantener una cola de tareas pendientes? (Es relativamente fácil para modificaciones puntuales, pero sigue siendo un reto para los flujos de trabajo continuos)
\end{itemize}
\end{warningbox}

\subsection{Resumen de la sección}

El modelo de Agent Manager exige que el desarrollador pase de ser «la persona que escribe código» a ser «la persona que gestiona a los Agents que escriben código». Las técnicas centrales incluyen archivos de comportamiento, hooks, commands y subagents. Los principios clave son: establecer barreras de protección, auditar las modificaciones, hacer coincidir la tarea con el modelo y hacer checkpoint con frecuencia.

\section{Conferencia invitada: la filosofía de diseño de Claude Code (Boris Cherny)}

\subsection{La programación está en un punto de inflexión}

Boris Cherny, el creador de Claude Code, compartió una idea clave: la productividad de los lenguajes de programación y de los IDE ha crecido siguiendo una curva exponencial, y la IA es precisamente la fuerza central que impulsa la ola más reciente de ese crecimiento.

\begin{knowledgebox}{El crecimiento exponencial de la productividad en la programación}
De Assembly a FORTRAN, de C a Python, de VS Code a Cursor/Claude Code: cada generación de lenguajes de programación e IDE ha logrado un incremento lineal de la productividad en una escala logarítmica. Las técnicas de verificación (verification) también han evolucionado en paralelo: de la depuración manual a los tipos estáticos, de las pruebas automatizadas al fuzz testing y al self-play impulsados por IA.
\end{knowledgebox}

\subsection{Los principios de diseño de Claude Code}

La filosofía de diseño central de Claude Code se resume en tres «omnipresencias»:

\begin{enumerate}
    \item \textbf{Terminal-native}: se ejecuta en la terminal, no como un complemento de IDE
    \item \textbf{Low-level model access}: proporciona acceso directo al modelo de bajo nivel
    \item \textbf{Infinitely hackable}: personalizable sin límites
\end{enumerate}

\subsection{Un solo Claude Code, múltiples formas}

Claude Code atiende distintos escenarios a través de múltiples formas:

\begin{itemize}
    \item \textbf{Terminal}: uso directo desde la línea de comandos
    \item \textbf{IDE}: integración en editores como VS Code
    \item \textbf{Web \& iOS}: acceso a través de \url{claude.ai/code}
    \item \textbf{GitHub App}: instalación mediante \texttt{/install-github-app}
    \item \textbf{SDK}: invocación programática, con salida en JSON y composición mediante tuberías
\end{itemize}

\subsection{Patrones típicos de uso de Claude Code}

\subsubsection{Preguntas y exploración sobre la base de código}

\begin{itemize}
    \item «How do I make a new ValidationTemplateFactory?»
    \item «Why did we fix issue \#18363 by adding the if/else in login.ts?»
    \item «Look at PR \#9383, then verify which app versions were impacted»
    \item «What did I ship last week?»
\end{itemize}

\subsubsection{Adaptar el flujo de trabajo a la tarea}

Distintas tareas se adaptan mejor a distintos patrones de flujo de trabajo:
\begin{itemize}
    \item \textbf{explore $\rightarrow$ plan $\rightarrow$ confirm $\rightarrow$ code $\rightarrow$ commit}: adecuado para bug fixes que requieren investigación previa
    \item \textbf{tests $\rightarrow$ commit $\rightarrow$ code $\rightarrow$ iterate $\rightarrow$ commit}: patrón TDD (desarrollo guiado por pruebas)
    \item \textbf{code $\rightarrow$ screenshot $\rightarrow$ iterate}: patrón de iteración visual para el desarrollo de UI
\end{itemize}

\subsubsection{Iteración rápida de prototipos}

Boris mostró un ejemplo vívido de iteración de prototipos: mediante instrucciones sucesivas en lenguaje natural, iteró rápidamente el diseño de la UI de una lista de tareas pendientes (TODO); pasando de una vista inicial del uso de herramientas, a un título en línea, a un panel independiente y, finalmente, a una visualización combinada con el spinner, cada iteración requirió solo un prompt.

\begin{importantbox}{Construir para el modelo de dentro de 6 meses}
La idea central que Boris subrayó: no diseñes tu flujo de trabajo para las capacidades del modelo de hoy, sino prepárate para las capacidades del modelo de \textbf{dentro de 6 meses}. Los modelos evolucionan rápidamente, y tus herramientas y flujos de trabajo también deben tener capacidad de evolucionar.
\end{importantbox}

\subsection{Restricciones organizacionales al gestionar múltiples Agents}

Otra conclusión implícita de la Week 4 es que, cuando los desarrolladores empiezan a gestionar varios Agents a la vez, la forma misma de gestionar proyectos también debe cambiar.

\begin{itemize}
    \item \textbf{Las tareas deben poder dividirse}: si los requisitos son ambiguos, es difícil delegarlos con seguridad a varios Agents en paralelo
    \item \textbf{Los límites deben ser claros}: es necesario especificar de antemano qué archivos, qué pasos de verificación y qué granularidad de commits corresponden a cada Agent
    \item \textbf{La revisión debe estandarizarse}: cuanto más se dependa de la ejecución paralela de Agents, más se necesita un mecanismo consistente de revisión de diffs y de reversión
\end{itemize}

\begin{warningbox}{El principal modo de falla del Agent Manager no es que el modelo no sea lo bastante inteligente, sino una división inadecuada de las tareas}
Muchos fracasos en la ejecución paralela no ocurren porque el Agent no pueda hacerlo, sino porque el humano no definió, antes de empezar, límites de responsabilidad ni criterios de aceptación claros. Dicho de otro modo, ser Agent Manager es, ante todo, una capacidad de gestión de proyectos, y solo en segundo lugar, una técnica de prompting.
\end{warningbox}

\subsection{Resumen de la sección}

La filosofía de diseño de Claude Code es terminal-native, acceso a modelos de bajo nivel y personalización sin límites. Cubre todo el SDLC, desde la exploración de la base de código hasta la gestión de CI/CD. La estrategia central de uso consiste en elegir el patrón de flujo de trabajo adecuado según el tipo de tarea, y prepararse con antelación para las capacidades futuras de los modelos.

\section{Síntesis y ampliación}

La semana 4 exploró los patrones de uso de los Coding Agents desde dos perspectivas complementarias: Mihail Eric explicó la metodología y las técnicas de gestión del Agent Manager, mientras que Boris Cherny compartió, desde la perspectiva del creador de Claude Code, la filosofía de diseño de herramientas y su experiencia práctica.

\subsection{Puntos clave}

\begin{itemize}
    \item El desarrollo de software está pasando de «personas que gestionan personas» a «personas que gestionan Agentes»: esta es la cuarta transformación de la forma organizacional
    \item Cuatro técnicas de orientación para Agentes: archivos de comportamiento, hooks, commands, subagents
    \item Prácticas clave: barreras de protección (tests + CI/CD), marcadores de auditoría, correspondencia de modelo, checkpoints frecuentes
    \item Filosofía de diseño de Claude Code: terminal-native, low-level access, infinitely hackable
    \item Elegir el flujo de trabajo según el tipo de tarea: explore-plan-code, TDD, iteración visual
    \item Prepararse para las capacidades del modelo dentro de 6 meses
\end{itemize}

\subsection{Lista de implementación}

\begin{enumerate}
    \item Construir commands, hooks y archivos de comportamiento reutilizables para tareas comunes
    \item Usar subagents para tareas paralelas con una división de trabajo clara, en lugar de apretar todo el trabajo en un solo contexto
    \item Tratar las pruebas, la CI y la revisión humana como límites estrictos del flujo de trabajo del Agente, y no como un remedio posterior
\end{enumerate}

\subsection{Tabla comparativa de división de roles}

\begin{table}[H]
\centering
\begin{tabular}{p{0.22\textwidth} p{0.28\textwidth} p{0.28\textwidth}}
\toprule
Rol & Responsabilidad principal & Forma de remediar en caso de falla \\
\midrule
Agent Manager humano & Definir objetivos, dividir tareas, establecer límites, tomar las decisiones finales & Redefinir las tareas, agregar restricciones, decidir si revertir o continuar \\
Agente ejecutor & Leer el contexto, modificar código, ejecutar pruebas, generar commits & Exponer los problemas mediante hooks / tests / review y reintentar \\
Sistema de revisión y CI & Proporcionar verificación objetiva, impedir que errores evidentes lleguen a la rama principal & Reducir riesgos ocultos mediante verificaciones estandarizadas y registros reproducibles \\
\bottomrule
\end{tabular}
\caption{Límites de responsabilidad entre humanos y máquinas en el modelo Agent Manager}
\end{table}

\subsection{Modelo de madurez de adopción en equipos}

Del contenido del curso se puede extraer una ruta de adoption muy práctica:

\begin{enumerate}
    \item \textbf{Etapa de experimentación individual}: ingenieros particulares usan Agentes para responder preguntas, generar código de andamiaje y hacer refactorizaciones locales
    \item \textbf{Etapa de estandarización en equipo}: comienzan a consolidarse commands, hooks, reglas de review y archivos de comportamiento unificados
    \item \textbf{Etapa de colaboración paralela}: los Agentes se usan en tareas multihilo con una división de trabajo clara, y los desarrolladores comienzan a asumir el rol de Agent Manager
    \item \textbf{Etapa de integración de sistemas}: el flujo de trabajo del Agente se integra profundamente con la CI, la revisión de código, las bases de conocimiento y las herramientas de gestión de proyectos
\end{enumerate}

\begin{knowledgebox}{La señal de que aumenta la madurez no es «se usan más Agentes», sino «las fallas son más controlables»}
Los equipos que en verdad alcanzan la etapa de madurez no son los que se ven más vistosos, sino aquellos en los que, cuando un Agente comete un error, las pruebas, los hooks, la revisión de código y los límites de responsabilidad lo detectan a tiempo con mayor facilidad. Por eso Boris y Mihail insisten repetidamente en los guardrails, y no solo en la velocidad.
\end{knowledgebox}

\subsection{Lecturas adicionales}

\begin{itemize}
    \item Sitio oficial de Claude Code: \url{https://claude.ai/code}
    \item Documentación del SDK de Claude Code: \url{https://docs.anthropic.com/en/docs/claude-code}
    \item Awesome Claude Agents: \url{https://github.com/vijaythecoder/awesome-claude-agents}
    \item SuperClaude Framework: \url{https://github.com/SuperClaude-Org/SuperClaude_Framework}
\end{itemize}

%% --- Fin del contenido --- %%

\end{document}
